\begin{proofbox}
	\textbf{\textcolor{CProof}{思路提示}}

	乘法公式的核心是条件概率定义的直接变形。对于多个事件，反复应用两事件乘法公式即可得到一般形式。

	\medskip
	\textbf{\textcolor{CProof}{正式推导}}
	\begin{description}[style=nextline, leftmargin=2.8em, labelsep=0.5em, itemsep=0.55em]
		\item[\textbf{步骤一（条件概率定义）：}] 设 $P(A)>0$。
			\[
				P(B|A) = \frac{P(AB)}{P(A)}
			\]
			\par\textbf{理由：}
			这是条件概率的定义式。两边乘以 $P(A)$，即得 $P(AB)=P(A)P(B|A)$。

		\item[\textbf{步骤二（三事件递推）：}] 对于事件 $A_1,A_2,A_3$，先合并 $A_1A_2$ 为一个事件。
			\[
				P(A_1A_2A_3) = P((A_1A_2)A_3) = P(A_1A_2)\,P(A_3|A_1A_2)
			\]
			\par\textbf{理由：}
			对 $A_1A_2$ 与 $A_3$ 应用两事件乘法公式。再用一次两事件公式于 $A_1A_2$：
			\[
				P(A_1A_2)=P(A_1)\,P(A_2|A_1)
			\]
			代入得 $P(A_1A_2A_3)=P(A_1)P(A_2|A_1)P(A_3|A_1A_2)$。

		\item[\textbf{步骤三（n事件归纳）：}] 假设对 $k$ 个事件公式成立，对于 $k+1$ 个事件 $A_1,\dots,A_{k+1}$，有
			\[
				P(A_1\cdots A_{k+1})=P(A_1\cdots A_k)\,P(A_{k+1}|A_1\cdots A_k)
			\]
			\par\textbf{理由：}
			将 $A_1\cdots A_k$ 视为一个事件，再次使用两事件公式。再由归纳假设，
			\[
				P(A_1\cdots A_k)=P(A_1)P(A_2|A_1)\cdots P(A_k|A_1\cdots A_{k-1})
			\]
			代入即得 $k+1$ 事件公式。故对任意正整数 $n$ 成立。
	\end{description}

	\medskip
	\textbf{\textcolor{CProof}{结论回扣}}

	\textbf{乘法公式将积事件概率转化为条件概率的连乘，是条件概率定义的直接推论，也是计算复杂事件概率的重要工具。}
\end{proofbox}
